\documentclass{article}
\usepackage[T1]{fontenc}
\usepackage{lmodern}
\usepackage{graphicx}
\usepackage{amsmath,amssymb}
\usepackage[a4paper,margin=2cm]{geometry}
\usepackage[absolute,overlay]{textpos}
\setlength{\TPHorizModule}{1mm}
\setlength{\TPVertModule}{1mm}
\usepackage[indonesian]{babel}
\usepackage{hyperref}
\setlength{\emergencystretch}{2em}
% unit: Halaman 87

\begin{document}

% === Halaman 87 (llm) ===

\section*{3 data Chunk}

Inilah bagian terpenting --- berisi sampel audio mentah.

Struktur:

\begin{tabular}{|l|l|l|}
\hline
\textbf{Field} & \textbf{Ukuran} & \textbf{Keterangan} \\ \hline
"data" & 4 byte & ID chunk \\ \hline
Subchunk2Size & 4 byte & Ukuran data audio \\ \hline
Data & n byte & Sampel audio \\ \hline
\end{tabular}

\vspace{10mm}

\subsection*{Representasi Data Audio}

Misalnya:
\begin{itemize}
    \item 16-bit
    \item Stereo
    \item 44100 Hz
\end{itemize}

Maka struktur sampel:

\begin{center}
\texttt{[L1][R1][L2][R2][L3][R3]...}
\end{center}

Setiap sampel:
\begin{itemize}
    \item 16-bit = 2 byte
    \item Stereo = 4 byte per frame
\end{itemize}

\end{document}
